\section{Does the arithmetic reproduce the excesses we have seen?}
\label{sec:excess}

A trials-factor calculation is easy to get wrong in a way that is invisible. Nothing in
Section~\ref{sec:budget} would look different if $\Nsig$ were off by an order of magnitude. There
is one external check available, though, and it is a strong one. If the number of independent
looks across the ATLAS search program is known even approximately, background-only statistics
predicts how many $3\sigma$ and $5\sigma$ local excesses should have appeared. That prediction can
be compared to the historical record.

\subsection{The prediction}

The relevant count is quasi-independent statistical tests, meaning mass-resolution elements plus
signal-region bins, across all ATLAS searches rather than bump hunts alone.

\begin{table}[htbp]
\centering
\small
\caption{Trials-factor anchors for the full ATLAS search program.}
\label{tab:excess_N}
\begin{tabular}{@{}l r >{\raggedright\arraybackslash}p{6.0cm}@{}}
\toprule
level & looks $N$ & basis \\
\midrule
bump-hunt budget (public channels) & $\sim3.7\times10^3$ & this study, Section~\ref{sec:budget} \\
two-body invariant-mass searches & $\sim2$--$3\times10^4$ & $>50\,000$ LHC-wide~\cite{statinterp}, ATLAS $\sim$ half \\
full ATLAS BSM program (incl.\ SRs) & $\sim5\times10^4$ & $300$--$500$ searches $\times\,\mathcal{O}(100)$ regions~\cite{diggingdeeper} \\
\bottomrule
\end{tabular}
\end{table}

With a one-sided Gaussian tail, $p(\ge3\sigma) = 1.35\times10^{-3}$ and $p(\ge5\sigma) =
2.87\times10^{-7}$, the expected counts are $N \cdot p$:

\begin{table}[htbp]
\centering
\caption{Expected upward fluctuations under background-only.}
\label{tab:excess_expected}
\begin{tabular}{lrr}
\toprule
$N$ & expected $\ge3\sigma$ & expected $\ge5\sigma$ (spurious) \\
\midrule
$10^4$ & 14 & 0.003 \\
$5\times10^4$ (central) & \textbf{67} & \textbf{0.014} \\
$10^5$ & 135 & 0.029 \\
resonance subset, $3.7\times10^3$ & 5.0 & $1.1\times10^{-3}$ \\
\bottomrule
\end{tabular}
\end{table}

\begin{figure}[htbp]
\centering
\includegraphics[width=0.85\textwidth]{excess_counting.png}
\caption{Expected number of upward fluctuations above $Z$ as a function of the number of
independent looks, with the anchors of Table~\ref{tab:excess_N} marked. The $\ge5\sigma$ curve
stays below unity across the entire defensible range of $N$, which is the quantitative content of
the $5\sigma$ convention~\cite{lyons2008}.}
\label{fig:excess}
\end{figure}

\subsection{The record}

\textbf{At $3\sigma$: order tens, essentially all faded.} The $750\GeV$ diphoton excess
($3.9\sigma$ local in 2015, $2.1\sigma$ global, gone with more data), the $2\TeV$ diboson $WZ$
excess at $8\TeV$, the $95\GeV$ and $151\GeV$ diphoton hints, the multilepton anomalies, a
$400\GeV$ $t\bar t$ excess, and the surviving $\sim3\sigma$ excesses found hiding inside published
SUSY signal regions~\cite{diggingdeeper}. That is the expected population, not a mystery.

\textbf{At $5\sigma$: every one is a genuine Standard-Model process.} The Higgs boson, four-top
production, $t\bar tZ$/$t\bar tW$/$t\bar tH$, vector-boson scattering, single top; the public
summary of ATLAS search results~\cite{atlas_public_searches} contains no BSM counterexample.
\textbf{Spurious BSM $\ge5\sigma$: exactly zero.} Four-top production is instructive rather than a
counterexample: it evolved from $3\sigma$ to $5\sigma$ and turned out to be real.

\subsection{Verdict}

\begin{table}[htbp]
\centering
\caption{Expected versus observed, for $N \sim 5\times10^4$.}
\label{tab:excess_verdict}
\begin{tabular}{lrl}
\toprule
 & expected & observed \\
\midrule
$\ge3\sigma$ local & $\sim15$--$135$ & order tens, all faded \\
$\ge5\sigma$ spurious & $\sim0.01$ & 0 \\
\bottomrule
\end{tabular}
\end{table}

The two agree. This matters for the note's central claim in a specific way: the same
look-elsewhere arithmetic that produces $\Zloc \approx 6.4$ in Section~\ref{sec:budget} also
post-dicts the observed excess population without adjustment. The framework is at least
self-consistent, and the parade of $3\sigma$ hints that come and go needs no explanation beyond $N
\times p$.

The corollary is uncomfortable. With more than $50\,000$ looks LHC-wide, a local $5\sigma$ in a
fully agnostic scan is worth only $1.5$ to $2.5\sigma$ globally, which is the basis of the
argument~\cite{statinterp} for raising the agnostic-scan threshold towards $7\sigma$. This note
reaches essentially the same number from the opposite direction for a combinatorial scan
(Section~\ref{sec:combinatorial}, $\Zloc = 6.98$).

Three caveats. $N$ spans a factor ten and is the dominant uncertainty here, although the
conclusion survives the whole range. The counts mix ATLAS-only and ATLAS$+$CMS excesses. And
``order tens observed'' counts \emph{notable} excesses; the raw number of individual $\ge3\sigma$
bins is much larger and goes unremarked, which is also consistent, since those are the expected
noise.
